
\section{Discussion and Conclusions}



We have presented the first approach that truly allows real-time, high-quality radiance field rendering, in a wide variety of scenes and capture styles, while requiring training times competitive with the fastest previous methods.

Our choice of a 3D Gaussian primitive preserves properties of volumetric rendering for optimization while directly allowing fast splat-based rasterization. Our work demonstrates that -- contrary to widely accepted opinion -- a continuous representation is \emph{not} strictly necessary to allow fast and high-quality radiance field training.

The majority ($\sim$80\%) of our training time is spent in Python code, since we built our solution in PyTorch to allow our method to be easily used by others. Only the rasterization routine is implemented as optimized CUDA kernels. We expect that porting the remaining optimization entirely to CUDA, as e.g., done in InstantNGP~\cite{mueller2022instant}, could enable significant further speedup for applications where performance is essential.

We also demonstrated the importance of building on real-time rendering principles, exploiting the power of the GPU and speed of software rasterization pipeline architecture. These design choices are the key to performance both for training and real-time rendering, providing a competitive edge in performance over previous volumetric ray-marching.

It would be interesting to see if our Gaussians can be used to perform mesh reconstructions of the captured scene. Aside from practical implications given the widespread use of meshes, this would allow us to better understand where our method stands exactly in the continuum between volumetric and surface representations.

In conclusion, we have presented the first real-time rendering solution for radiance fields, with rendering quality that matches the best expensive previous methods, with training times competitive with the fastest existing solutions.

